\section{Discussion, Limitations, and Future Work}
\label{sec:discussion}

\subsection{Experimental Protocol and Sample Accounting}
\label{subsec:accounting}

Different subsets of videos appear across our experiments, which we clarify here.
Table~\ref{tab:protocol} summarises inclusion criteria, subset sizes, and the claim each subset supports.

\begin{table}[h]
\centering
\small
\caption{Protocol and sample accounting across experiments. ``Valid'' = videos passing $\mathrm{JF}_\mathrm{clean} \geq 0.30$ and successful codec round-trip.}
\label{tab:protocol}
\begin{tabular}{lrlll}
\toprule
\textbf{Experiment} & \textbf{N} & \textbf{Dataset} & \textbf{Codec} & \textbf{Claim} \\
\midrule
Main DAVIS paired       & 36  & DAVIS (90 total) & H.264 CRF=23 & AdvOpt vs.\ baseline gain \\
Global-fixed DAVIS      & 27  & DAVIS            & H.264 CRF=23 & Collapse to universal point \\
Proxy validation        & 88  & DAVIS            & H.264 CRF=23 & Within-video proxy $r{=}+0.811$ \\
Oracle-gap              & 88  & DAVIS            & H.264 CRF=23 & 98\% proxy-oracle match \\
Utility (SSIM/YOLO)     & 20  & DAVIS            & H.264 CRF=23 & Privacy-utility tradeoff \\
HEVC generalization     & 17  & DAVIS            & HEVC CRF=23  & Codec generalization \\
XMem cross-tracker      & 19  & DAVIS            & H.264 CRF=23 & Tracker generalization \\
YT-VOS full sweep       & 497 & YT-VOS (507 valid)& H.264 CRF=23 & Dataset generalization \\
YT-VOS op.\ point       & 50  & YT-VOS           & — (opt.\ only)& $\alpha$ convergence at ssim\_floor=0.90 \\
YT-VOS fair comparison  & 50  & YT-VOS           & H.264 CRF=23 & idea1($\alpha{=}0.93$) vs.\ AdvOpt at matched quality \\
\bottomrule
\end{tabular}
\end{table}

The DAVIS main sweep uses $n=36$ of 90 videos; the full 90-video sweep is ongoing.
The YT-VOS fair comparison ($n=50$, ssim\_floor=0.90) directly compares idea1($\alpha{=}0.93$) vs.\ AdvOpt under the same codec protocol:
AdvOpt achieves $16.45\,\mathrm{pp}$ vs.\ idea1's $17.27\,\mathrm{pp}$ (mean gain $= -0.82\,\mathrm{pp}$, $t{=}{-}2.08$, $n{=}50$);
$94\%$ of gains fall within $\pm 5\,\mathrm{pp}$, confirming near-equivalence for the majority of videos.
The $-0.82\,\mathrm{pp}$ difference traces entirely to $6/50$ videos where the MSE surrogate over-penalises and the optimizer backs off unnecessarily (surrogate brittleness; see Section~\ref{subsec:soft_constraint}).

\subsection{Operating-Point Discovery vs.\ Per-Video Adaptation}

The central finding---that a global-fixed setting $(r=24,\alpha=0.93)$ matches per-video AdvOpt on DAVIS ($n=27$, $+0.3\,\mathrm{pp}$, $t=0.31$)---should be interpreted carefully.
It does not mean AdvOpt is unnecessary. It means AdvOpt \emph{discovered} that $\alpha\approx 0.93$ is the near-universal codec-robust operating point, which was not known \emph{a priori}.
The intuitive starting point is $\alpha=0.80$ (idea1), which leaves $\sim 33\,\mathrm{pp}$ of privacy gain on the table.
The proxy (within-video $r=+0.811$, oracle gap 98\%, both on DAVIS) validates that AdvOpt reliably finds this optimal point under the quality constraint, rather than by exhaustive search.
We note that proxy validation is conducted on DAVIS; transfer to YT-VOS and HEVC is supported by the convergence of AdvOpt's selected $\alpha$ to the same operating point ($94\%$ of YT-VOS videos converge to $\alpha\in[0.92,0.93)$ under ssim\_floor=0.90), but within-video proxy correlation on YT-VOS is not yet available and is left for future work.

The YT-VOS fair comparison ($n=50$) directly tests this: idea1($\alpha{=}0.93$) vs.\ AdvOpt($\mathrm{ssim\_floor}{=}0.90$) on the same 50 videos with the same codec protocol.
AdvOpt converges to $\alpha{>}0.92$ for $88\%$ of videos, matching idea1 closely (within $\pm 5\,\mathrm{pp}$ for $94\%$ of pairs).
For $12\%$ of videos (6/50), the MSE surrogate over-penalises: AdvOpt backs off to $\alpha{<}0.92$ even though idea1's actual SSIM $\geq 0.96$ for all 6 of these videos, indicating false constraint activation rather than a genuine quality risk.
The net effect: AdvOpt achieves $-0.82\,\mathrm{pp}$ less privacy degradation on average ($t{=}{-}2.08$, $n{=}50$), driven entirely by these surrogate-brittle cases.
Both methods have the same SSIM violation rate ($3/50 = 6\%$ below $0.90$).
The correct deployment strategy is therefore: use AdvOpt offline to discover the global operating point ($r^*{=}24$, $\alpha^*{\approx}0.93$), then apply this fixed setting globally.

\subsection{Soft Surrogate Constraint: Scope and Violations}

AdvOpt enforces quality via a soft surrogate: a differentiable MSE penalty acting as a proxy for $\mathrm{SSIM} \approx 1 - 20 \cdot \mathrm{MSE}$.
This is not a hard post-codec constraint. The surrogate is tight in the low-distortion regime but can underestimate the true SSIM impact on high-contrast scenes.
We observe a $6\%$ soft-violation rate (realized SSIM $< 0.90$) on YT-VOS for AdvOpt, and $2\%$ for the fixed baseline ($\alpha=0.80$).

We report this honestly: the method maintains SSIM $\geq 0.90$ for $94\%$ of YT-VOS videos and for all tested DAVIS videos with the global-fixed setting.
The surrogate brittleness is documented in our ssim\_floor ablation (Section~\ref{subsec:accounting}): at ssim\_floor=0.92 with a strong baseline ($\alpha=0.93$), the optimizer collapses to very low $\alpha$ because the constraint binds too aggressively.
The correct operating regime is ssim\_floor $\in [0.88, 0.90]$, where the surrogate provides reliable guidance.
Future work should replace the MSE-SSIM surrogate with the true SSIM gradient or a learned quality predictor.

\subsection{Privacy-Utility Tradeoff}

AdvOpt achieves stronger privacy ($+57.7\,\mathrm{pp}$, SSIM $\approx 0.928$) than the fixed baseline ($+24.7\,\mathrm{pp}$, SSIM $\approx 0.953$) at a measurable downstream cost: YOLO detection recall drops from $77.5\%$ (baseline) to $55.2\%$ (AdvOpt).
This is a deliberate privacy-utility tradeoff, not a side effect.
The appropriate use case is a dataset publisher who prioritises strong, codec-robust privacy guarantees and can accept reduced downstream utility for third-party annotators.

\subsection{Scope of Generalization}

Cross-codec generalization (HEVC: $+28.5\,\mathrm{pp}$, $t=4.20$) confirms that the operating point is robust to codec-family change.
Cross-tracker transfer to XMem ($17.6\,\mathrm{pp}$) is present but attenuated (vs.\ $\sim 54\,\mathrm{pp}$ for SAM2), consistent with XMem's memory-based matching being less reliant on boundary features than SAM2's contour-heavy encoder.
The method is not tracker-agnostic: its effectiveness depends on the target tracker's boundary-feature dependency.

\subsection{Limitations}

\paragraph{Proxy scope.} The proxy and oracle-gap experiments are validated on DAVIS ($n=88$). YT-VOS proxy validation (within-video correlation) is not yet available. The convergence of AdvOpt's selected $\alpha$ to $[0.92, 0.93)$ on $94\%$ of YT-VOS videos provides indirect evidence of proxy reliability, but a direct within-video correlation study on YT-VOS is needed to confirm generalization. We therefore restrict proxy-validated claims to DAVIS in the main results, with YT-VOS and HEVC presented as generalization evidence rather than proxy validation.

\paragraph{Single backbone.} All experiments use SAM2.1 hiera\_tiny. Larger SAM2 variants (Small, Large) may differ in boundary-feature reliance.

\paragraph{Mask quality.} AdvOpt uses the first annotated frame only. For imperfect or missed masks, the selected parameters may be suboptimal.

\paragraph{Differentiable codec.} A true end-to-end codec-in-loop optimizer could potentially discover stronger operating points than the boundary gradient proxy. This is the most important direction for future work.
